% !TeX root = ../papers/nc-p3.tex
\section{平行な二重色と四つの非欠損長方形の全被覆}
\label{sec:full-thin-parallel-cover}
\begin{theorem}\label{thm:full-thin-parallel-cover}
四元群が四位置で正則表現型に作用し、二つの隣接切断の色階数がともに $(2,2,1,1)$ の
並べ替えで、二つの厚い色の差が一致するとする。四つの薄い階数1長方形の全ての行・列が
非零なら、全ての標準Koszul左右AS窓は、文字の付け替えと始点の移動を含めて、
次の三族の閉包で覆われる。
\begin{enumerate}
\item 定理\ref{thm:klein-elliptic-completion}の五次元族。
\item $\lambda=1,B=A$ の両対称な四次元交換表示。
\item その非自明コサイクル捻り。
\end{enumerate}
これは四つの薄い長方形の係数だけに非消滅を課した被覆であり、
十六個の厚い係数のうちいくつが消えてもよい。
異なる厚色方向では定理\ref{thm:klein-transverse-dense}の二変数表示が適用されるため、
二重色の両方向について、薄い欠損のない場合に正則候補を含む族が尽くされた。
\end{theorem}
\begin{proof}[薄い係数の商と切断の入替え]
節\ref{sec:klein-elliptic-completion}の規格化を使う。薄い48語は二文字ずつが両半分にある語で、
係数は $z,\lambda$ のLaurent単項式である。各語の四位置での文字指数行列を $E$、
その転置の整数核の基底を行とする行列を $L$、薄い係数の二指数の行列を $F$ とすれば、
$LF$ の整数行格子は $\mathbb Z^2$ そのものである。従って対角文字倍率を除いた薄い情報は
この二単元で正確に表される。共通の線型文字置換と位置別平行移動で両厚色集合 $\{0,1\}$ を
保存する32通りを考える。線型部は恒等または2と3の交換、四移動量は総和零で全上位ビットが等しい。
置換・始点移動後の指数行列を $F'$ として $L F'=LF M$ を解くと、
$(t,v)=(\lambda,z\lambda)$ に対する作用は次の通りである。
始点固定では $(t,v)$ の交換と同時逆数が生成する四元群であり、
一位置移動を許せば独立な逆数と交換が生成する位数8の全群となる。
\path{computations/verify_full_thin_exchange_cover.py} は位置移動の偶奇別に全候補を検算する。
これらは実際の対角基底変更を伴う。代数閉体上でトーラス写像の像はその整数文字関係で決まるからである。

$x=t+t^{-1},y=v+v^{-1}$ と書くと、始点を指定した薄い表示の商は
\[
 S=x+y,\quad P=xy,\quad D=(t-t^{-1})(v-v^{-1}),\qquad
 D^2=(P+4)^2-4S^2 .
\]
実際に同時逆数の不変環は $x,y,D$ で生成され、交換は $x,y$ を交換して $D$ を固定する。
一位置移動は $D$ の符号を変える。この商は指定した文字作用についてのものである。
同じ二点が $\{z,z\lambda^2\}$ について順序交換と独立な逆数を受ける。
数 $z$ と $z\lambda^2$ の一方が1なら、切断の向きを選んで $z=1$ にできる。
他方も1なら定理\ref{thm:full-thin-double-exchange}による。第二型は五次元族へ、
第一型の正負は定理\ref{thm:klein-bisymmetric-irreducible}により対応する交換表示の閉包に入る。
他方が1でなければ定理\ref{thm:klein-mixed-exchange-resolution}による。

それ以外は両切断の交差ブロックが全て可逆であり、各厚い角も階数2条件から可逆である。
$z,z\lambda^2$ の一方が $-1$ とも異なれば、そちらを $z$ にして
定理\ref{thm:klein-all-pairings}を適用できる。四つの例外表示の二つの正表示は
$(t,v)$ の交換で相互に移る。具体的に $v=1$ の一般点は $t'=1$ へ移り、
正交換の一般点は両対称なので負交換の一般点とは反対階数により異なる。
五次元主枝は $A,z,\lambda$ の一般値の上の既約な解であり、
いずれの付け替えもこの次元の唯一の枝を保つ。
従って付け替え後にも両対称例外枝の稠密開部は主枝の外にある。
同じことを捻りの後にも使える（薄い符号は $\lambda\mapsto-\lambda$）。
四例外表示は各々既約なので、一般点での帰着はその閉包全体の帰着を与える。
\end{proof}

\begin{proof}[最後の繊維 $z=z\lambda^2=-1$]
コサイクル捻りで $\lambda=1$ とする。
$A=\left(\begin{smallmatrix}a&b\\c&d\end{smallmatrix}\right)$ に対し
$(p,q,r,w)=(a+d,a-d,b+c,b-c)$、$B$ の対応する四数を $(P,Q,R,W)$ と書く。
節\ref{sec:klein-elliptic-completion}の角 $C,D$ を使い、
$u=\det C/\det A$, $v=\det D/\det A$, $t=\det B/\det A=u+v-1$ とする。
いずれも非零である。$X=\left(\begin{smallmatrix}1&1\\1&-1\end{smallmatrix}\right)$,
$Y=\left(\begin{smallmatrix}1&1\\-1&1\end{smallmatrix}\right)$ とすると
\[
 Y\operatorname{adj}(A) X=
 \begin{pmatrix}p-r&-q+w\\q+w&-p-r\end{pmatrix}.
\]
$C$ は $(p,q,R,W)$、$D$ は $(P,Q,r,w)$ で与えられる。
従って二つの切断の条件はちょうど
\[
 R=ur,\quad Q=vq,\quad (u-1)(p,P)=0,\quad (v-1)(w,W)=0,
 \quad (u-1)(v-1)(q,r)=0
\]
である。最後の式は $uv-t=(u-1)(v-1)$ から来る。
$u,v$ がともに1と異なる場合は $A=0$ となるので解がない。

$u=v=1$ なら $R=r,Q=q$, $P=\pm p,W=\pm w$ で尽きる。
この四表示は $B=A,A^{\mathsf t},
\left(\begin{smallmatrix}-d&b\\c&-a\end{smallmatrix}\right),
\left(\begin{smallmatrix}-d&c\\b&-a\end{smallmatrix}\right)$ の順である。
第一は正交換、第三は $b\ne c$ なら節\ref{sec:klein-doublet-boundary}の
第二補完を $z\to-1$ とした極限そのものであり、$b=c$ にも特殊化する。
第二は一位置移動を含む下記の第一変換で第三へ、第四は捻りを含む第二変換で第一へ移る。
前者の第二補完枝は $\lambda=1$ で五次元族に含まれた有理切断である。
実際八角を $(a,b,c,d,e,f,g,h)$ と並べると二変換は
\[
 (g,a,-d,-f,c,e,-h,-b),\qquad (b,a,-d,-c,g,-h,e,-f)
\]
である。いずれも後二文字に1を足す。第一は一位置の巡回を、第二は非自明捻りを伴い、
薄い符号を合わせる文字倍率だけで上式を得る。薄い語上の整数Smith方程式と八角で検算でき、
高位角はSchur補完で一意である。

$u=1,v\ne1$ なら $w=W=0,R=r$ なので二つの非対角はすべて $b$ で一致する。
$A,B$ は対称であり、残りは
\[
 P^2=v p^2+(1-v)r^2+v(v-1)q^2,\qquad Q=vq .
\]
これがこの側の全解である。同じ非対角を持つため偶交換そのものが成立する
（両切断で $C=A,D=B$ が成立する）。相異なる補完は全て既出の第二補完の閉包に入る。
直接には $e=(kd-b^2)/(d-k)$, $h=(ka-b^2)/(a-k)$ と取って $k$ を動かせばよい。
この式は $(e-h)(ad-b^2)=(a-d)(eh-b^2)$ と同値な一般表示であり、
上記の方程式は $v$ に関する二次式の平方根のモデルで既約である。
実際右辺は $p,q,r$ の階数3の二次形式で、$v(v-1)\ne0$ で平方にならない。
よって特殊値も閉包に入る。$v=1,u\ne1$ では $p=P=0,Q=q$ であり、
上の第一変換で対称な二角・共通非対角の側に戻る。
これが最後の繊維の全被覆を与える。
\end{proof}

\begin{corollary}\label{cor:minus-one-second-completion-periods}
第三表示 $B=\left(\begin{smallmatrix}-d&b\\c&-a\end{smallmatrix}\right)$ と
$u=1,v\ne1$ の表示はともに反対階数10となる稠密開集合を持つ。
両者の各々の稠密開集合の全閉点は標準Koszul左右AS正則である。
第三表示の幾何学的一般点の最小周期は4、射影安定群の位数は8である。
$u=1,v\ne1$ 表示の反対階数10の正則点は全て最小周期2である。
\end{corollary}
\begin{proof}
前者で $(a,b,c,d)=(2,3,5,7)$、後者で $(a,b,d,e,h)=(2,3,7,5,1)$ とする。
両点で反対階数10、全始点で五次までの経路次元は標準値、四次中心の色別次元は
$(3,2,0,0)$。有理中心四式で割った商の素数1009での次元は全始点で次数13まで
$(1+t+t^2+t^3)^4$ の係数である。定理\ref{thm:klein-generic-central-parameters}の
第二補完枝の一般中心次元に等しいので、この二点の周囲で中心基底が延び、
有限個の商次元上界と全次数のHilbert下界を比較する自由加群の証明が適用できる。

前者の標本の素数109での四つの配置は各々相異なる二十点で、
射影写像数の四行は $(8,0,8,0),(0,8,0,0),(8,0,8,0),(0,0,0,8)$ である。
三点直線を許して、各点を通る三点直線上の三つ組と大きさ別の多点平面を数え、残る枠の像を全列挙した。
被約二十点の固有な枠比較でこの上界を一般点に使える。位置1と3の非同値性で周期2が除かれる。
位置0の射影安定群が窓の安定群の上界になることは前節の有限グラフの場合と同じである。
四元群に加えて位置別移動 $(1,0,1,0)$ が文字倍率補正により窓を保存する。
実際その八角の変換は共通倍率を除き $(h,-f,-g,e,d,-b,-c,a)$ であり、
第三表示では元の八角の $-1$ 倍になる。よって位数8が達成される。
後者の表示では二位置移動が文字倍率を除き $(A,B)\mapsto(A^{\mathsf t},B^{\mathsf t})$ を与え、
両角が対称なので周期2を持つ。周期1なら偶交換を一位置運んで奇交換も得る。
反対対は互いに素であるため両同定を同時に使えて両対称となるが、
定理\ref{thm:bisymmetric-six-or-eight}に反する。
\end{proof}

\begin{corollary}\label{cor:mixed-ordinary-open-center}
定理\ref{thm:klein-mixed-exchange-resolution}の四次元表示でも、同節の二つの有理点を
含む一つの稠密開集合の全閉点が四中心四次式上自由であり、標準Koszul左右AS正則である。
\end{corollary}
\begin{proof}
五次元楕円族は定理\ref{thm:full-fifteen-component}の一意な十五次元成分に含まれ、
従って今回の混合表示もそこに含まれる。この成分の一般の四次中心は五次元である。
二つの有理点で全始点の長さ4,5の次元は標準値に等しく、経路束はその近くで局所自由である。
中心方程式の核次元はそこで5以上で、二点ではちょうど5である。よって中心はその近くで部分束となり、
認証に使った四元を局所切断として延ばせる。中心商の長さ13までの次元上界は開条件なので、
二点それぞれの近傍の和で自由加群と全次数正則性の証明が通用する。
\end{proof}
